\section{Vision arquitectonica}

\subsection{Descripcion}

La solucion se documenta como un sistema web dividido en tres repositorios: frontend, backend y documentacion. El dominio del juego es el centro del sistema y condiciona el diseno de las demas capas.

\subsection{Alcance}

Esta vista describe componentes, limites y principios de separacion. No reemplaza la documentacion detallada de backend ni frontend.

\subsection{Definiciones}

\begin{itemize}
  \item \textbf{Repositorio frontend}: aplicacion web construida en Next.js.
  \item \textbf{Repositorio backend}: API y logica de negocio construida en Node.js y Express.
  \item \textbf{Repositorio documental}: fuente de verdad de reglas, dominio y criterios del sistema.
  \item \textbf{Modular monolith}: backend dividido por modulos de dominio con despliegue unico.
\end{itemize}

\subsection{Reglas}

Vista de alto nivel:

\begin{itemize}
  \item Frontend consume al backend exclusivamente mediante contratos de API.
  \item Backend concentra autenticacion, autorizacion, reglas de negocio, integraciones y persistencia.
  \item PostgreSQL almacena usuarios, comprobantes, partidos, picks, resultados, puntajes y auditoria.
  \item Una API externa provee resultados, con capacidad de override manual.
  \item La documentacion define reglas antes de que se cierren interfaces de implementacion.
\end{itemize}

Principios estructurales:

\begin{itemize}
  \item El dominio no debe depender de detalles de UI.
  \item La logica de negocio no debe vivir en controllers.
  \item La validacion de entradas ocurre en la frontera de entrada.
  \item Las acciones criticas deben quedar auditadas.
  \item El sistema debe poder recalcular sin romper consistencia.
\end{itemize}

\subsection{Casos borde}

\begin{itemize}
  \item Si una necesidad de implementacion entra en conflicto con una regla del dominio, la regla se revisa documentalmente antes de ajustarse en codigo.
  \item Si el proveedor externo de resultados cambia, la arquitectura debe preservar el contrato interno del dominio.
\end{itemize}

\subsection{Invariantes}

\begin{itemize}
  \item El frontend no accede directamente a la base de datos.
  \item El backend es responsable de enforcement de locks y scoring.
  \item La documentacion es el origen de definiciones funcionales.
\end{itemize}

\subsection{Preguntas abiertas}

\begin{itemize}
  \item Ninguna por ahora.
\end{itemize}

